\documentclass{template/ieeetj}

\makeatletter
\def\input@path{{template/}}
\makeatother

\usepackage{template/paper-template}

\title{{{TITLE}}}
\author{{{AUTHORS_LATEX}}}
\affil{Independent researcher}

\authornote{Demonstration manuscript for a structured critical reading.}

\begin{document}
\receiveddate{XX Month, XXXX}
\reviseddate{XX Month, XXXX}
\accepteddate{XX Month, XXXX}
\publisheddate{XX Month, XXXX}
\currentdate{\today}
\doiinfo{DOI pending}

\maketitle

\begin{abstract}
This demonstration paper presents a compact method for critically reading a
scientific article. The method separates claims, evidence, assumptions,
validity threats, and alternative explanations so that criticism remains
traceable to the source rather than relying on impressionistic judgment.
\end{abstract}

\section{Introduction}

A critical reading should reconstruct an article's argument before evaluating
it. The reader first identifies the research question and principal claim,
then determines which observations and analytical steps are intended to
support that claim. This reconstruction prevents disagreement with a topic
from being mistaken for evidence against the reported study.

\section{Method}

The review records the article's stated contribution, study design, data
source, comparison strategy, uncertainty treatment, and reproducibility
artifacts. It then tests whether each inferential step is supported, whether
important assumptions are explicit, and whether alternative explanations
remain credible.

\begin{enumpar}
  \item identify the article's principal claim
  \item trace the evidence supporting that claim
  \item test the assumptions connecting evidence to inference
\end{enumpar}

\section{Discussion}

The resulting critique should distinguish a demonstrated defect from a
limitation that merely narrows generalization. A useful review explains the
consequence of each issue and identifies the additional evidence or analysis
that would resolve it.

\section{Conclusion}

Structured critical reading improves rigor by making both praise and
criticism accountable to identifiable claims, methods, and evidence.

% Add \bibliographystyle{IEEEtran} and \bibliography{refs/refs} after
% introducing the manuscript's first citation.

\end{document}
